\documentclass[11pt]{article}
\usepackage{campaign}

\title{Executive Summary: Capacity-Based MAiD Reform}
\author{Justin Bogner}
\date{2026}

\begin{document}

\maketitle

\begin{abstract}
\noindent This policy brief outlines a fundamental structural reform to Medical Aid in Dying (MAiD) eligibility: replacing subjective, suffering-based criteria (such as terminal prognoses) with a rigorous, objective standard of decisional capacity. 
\end{abstract}

\section*{The Problem: Incoherent Eligibility}
Across North America, MAiD eligibility relies on clinical assessments of suffering (e.g., a 6-month terminal prognosis or ``grievous and irremediable suffering''). This framework is flawed in two key ways:
\begin{enumerate}
    \item \textbf{Philosophically Incoherent:} Patients hold the absolute right to refuse life-sustaining treatment based purely on \textit{decisional capacity}. Restricting active MAiD to specific diagnoses creates an asymmetry in bodily autonomy.
    \item \textbf{Subjective and Unequal:} Suffering cannot be objectively measured, forcing physicians to make subjective moral judgments about whose life is worth living.
\end{enumerate}

\section*{The Solution: A Decisional Capacity Standard}
The state must replace all suffering-based criteria with a single, rigorous standard: \textbf{decisional capacity}. A person who demonstrates the capacity to make an informed, durable decision regarding their own death should have access to the means to safely exit, regardless of medical diagnosis or disability status.

\section*{Four Structural Commitments}
The proposed model legislation rests on four interlocking pillars:

\subsection*{1. The Sovereignty Model}
Bodily self-determination is a fundamental right. Because the state maintains a monopoly on the pharmacological means of a peaceful exit, it holds a corresponding duty to provide access to any citizen of sound mind.

\subsection*{2. Two-Stage Capacity Assessment}
To replace the terminal illness restriction, the framework institutes rigorous procedural safeguards:
\begin{itemize}
    \item \textbf{Stage One:} Initial capacity assessment by an attending physician.
    \item \textbf{Stage Two:} A mandatory waiting period followed by an independent psychiatric review. Depression does not automatically disqualify a patient unless it actively impairs their ability to make a durable decision.
\end{itemize}

\subsection*{3. The Forcing Function (Accountability)}
To address concerns that systemic poverty or inadequate care will drive MAiD requests, the framework introduces the structural accountability mechanism:
\begin{itemize}
    \item Every MAiD request triggers \textbf{mandatory socioeconomic driver reporting}.
    \item This generates political accountability for systemic failures without paternalistically denying the individual their right to exit.
\end{itemize}

\subsection*{4. Fiscal Transparency}
While full fiscal impact modeling is provided in the Master Policy, strict ethical guardrails dictate that fiscal savings must \textit{never} function as a rationale for legalization.

\vspace{1em}
\noindent \textbf{Full Policy Packet \& Draft Legislation available at:} \url{https://autonomyandcapacity.org} \\
\noindent \textbf{Contact:} Justin Bogner \textbar{} (424) 666-1618 \textbar{} justin@cambrianminds.com \textbar{} 934 Iowa Street, Huntington, IN 46750

\end{document}
